\documentclass[12pt,a4paper]{article}

\usepackage[utf8]{inputenc}
\usepackage[T1]{fontenc}
\usepackage[spanish,es-nodecimaldot]{babel}
\usepackage{lmodern}
\usepackage{geometry}
\usepackage{setspace}
\usepackage{parskip}
\usepackage{microtype}
\usepackage{graphicx}
\usepackage{booktabs}
\usepackage{longtable}
\usepackage{array}
\usepackage{tabularx}
\usepackage{enumitem}
\usepackage{xcolor}
\usepackage{hyperref}
\usepackage{listings}
\usepackage{amsmath}
\usepackage{amssymb}
\usepackage{tikz}
\usetikzlibrary{arrows.meta,positioning,shapes.geometric}

\geometry{top=2.4cm,bottom=2.4cm,left=2.6cm,right=2.6cm}
\onehalfspacing
\hypersetup{
    colorlinks=true,
    linkcolor=blue!45!black,
    urlcolor=blue!60!black,
    citecolor=blue!45!black,
    pdftitle={RecruitAI: Agente Autónomo de Reclutamiento Inteligente},
    pdfauthor={Vane, Xio y Sandra}
}

\definecolor{primary}{HTML}{174A7E}
\definecolor{secondary}{HTML}{2D7D9A}
\definecolor{lightblue}{HTML}{EAF3F8}
\definecolor{codebg}{HTML}{F5F7FA}
\definecolor{codeborder}{HTML}{D9E1E8}
\definecolor{darkgray}{HTML}{3C4858}

\lstdefinestyle{professional}{
    backgroundcolor=\color{codebg},
    basicstyle=\ttfamily\small,
    breaklines=true,
    frame=single,
    rulecolor=\color{codeborder},
    numbers=left,
    numberstyle=\tiny\color{darkgray},
    keywordstyle=\color{primary}\bfseries,
    commentstyle=\color{green!45!black},
    stringstyle=\color{orange!55!black},
    showstringspaces=false,
    tabsize=4
}
\lstset{style=professional}

\newcommand{\projectname}{RecruitAI}
\newcommand{\projecttitle}{Agente Autónomo de Reclutamiento Inteligente}
\newcommand{\editable}[1]{\textcolor{darkgray}{\textit{#1}}}

\begin{document}

% ----------------------------------------------------------
% Portada
% ----------------------------------------------------------
\begin{titlepage}
    \centering
    \vspace*{1.2cm}

    {\Large \editable{NOMBRE DE LA UNIVERSIDAD}\par}
    \vspace{0.35cm}
    {\large \editable{Facultad o Escuela Profesional}\par}
    \vspace{1.8cm}

    {\color{primary}\rule{\textwidth}{1.2pt}\par}
    \vspace{0.55cm}
    {\Huge\bfseries \projectname\par}
    \vspace{0.3cm}
    {\LARGE\bfseries \projecttitle\par}
    \vspace{0.55cm}
    {\color{primary}\rule{\textwidth}{1.2pt}\par}

    \vspace{1.5cm}
    {\Large Informe técnico profesional\par}
    \vspace{0.5cm}
    {\large Sistema web con análisis explicable de CV, cola persistente de eventos y worker autónomo\par}

    \vfill

    \begin{tabular}{>{\bfseries}rl}
        Curso: & \editable{Nombre del curso} \\
        Docente: & \editable{Nombre del docente} \\
        Integrantes: & Vane --- Interfaz, arquitectura y flujo del agente \\
                     & Xio --- Extracción, NLP y compatibilidad \\
                     & Sandra --- Ranking, correos, reportes e historial \\
        Fecha: & \today
    \end{tabular}

    \vfill
    {\large Lima, Perú\par}
\end{titlepage}

\pagenumbering{roman}
\tableofcontents
\newpage
\listoftables
\newpage
\pagenumbering{arabic}

% ----------------------------------------------------------
\section{Resumen ejecutivo}
% ----------------------------------------------------------

\projectname{} es una plataforma web de reclutamiento desarrollada con Flask y SQLite. El sistema permite registrar vacantes, administrar candidatos, cargar currículums vitae (CV), analizar perfiles, calcular compatibilidad, generar rankings explicables, producir reportes en PDF y Excel, preparar invitaciones a entrevistas y enviarlas mediante un servidor SMTP.

La característica principal del trabajo es la incorporación de un \textbf{agente autónomo}. La aplicación web y el motor de análisis se ejecutan como procesos independientes. Flask funciona como panel de operación y monitoreo, mientras que un worker consulta una cola persistente de eventos almacenada en SQLite. De esta forma, el análisis no depende de que una persona mantenga abierta una página o presione un botón para cada CV.

El sistema implementa separación de roles entre el personal de Recursos Humanos (RRHH) y los candidatos. Asimismo, integra validación de archivos, protección CSRF, almacenamiento seguro de contraseñas, control de acceso por rol, historial de acciones, reintentos automáticos y recuperación de eventos interrumpidos.

\textbf{Palabras clave:} agente autónomo, reclutamiento inteligente, Flask, SQLite, análisis de CV, procesamiento asíncrono, SMTP, ranking explicable.

% ----------------------------------------------------------
\section{Introducción}
% ----------------------------------------------------------

Los procesos de selección suelen requerir la revisión manual de numerosos CV, la comparación repetitiva de requisitos y la coordinación de entrevistas. Estas actividades consumen tiempo y dificultan mantener criterios uniformes. \projectname{} propone una solución local y auditable que automatiza las tareas repetitivas sin eliminar la supervisión humana en decisiones sensibles.

El proyecto utiliza técnicas de procesamiento de texto orientadas a extracción de información y reglas ponderadas de compatibilidad. No depende de servicios externos de inteligencia artificial para razonar sobre los perfiles. Esta decisión facilita su ejecución local, reduce costos y permite comprender por qué se asignó una puntuación determinada.

La solución también incorpora autonomía operativa. El agente observa eventos del sistema y archivos nuevos, ejecuta análisis en segundo plano, actualiza rankings, genera reportes y prepara acciones posteriores. RRHH conserva el control final antes de enviar las invitaciones por correo.

% ----------------------------------------------------------
\section{Planteamiento del problema}
% ----------------------------------------------------------

\subsection{Situación identificada}

En un proceso convencional, RRHH debe realizar manualmente las siguientes actividades:

\begin{itemize}[leftmargin=*]
    \item recibir CV desde diferentes canales;
    \item revisar datos personales, experiencia, estudios y habilidades;
    \item comparar cada perfil con los requisitos de una vacante;
    \item priorizar candidatos;
    \item generar reportes;
    \item preparar mensajes de entrevista;
    \item registrar el historial de acciones.
\end{itemize}

Cuando aumenta el número de postulantes, estas tareas se vuelven lentas y propensas a inconsistencias. Además, una aplicación web tradicional puede limitarse a ejecutar lógica únicamente cuando recibe una solicitud HTTP. En ese caso existe automatización, pero no autonomía real.

\subsection{Pregunta de trabajo}

\begin{quote}
¿Cómo diseñar un sistema web de reclutamiento que analice CV de manera explicable, mantenga separados los permisos de RRHH y candidatos, y ejecute tareas automáticamente fuera del ciclo de peticiones HTTP?
\end{quote}

% ----------------------------------------------------------
\section{Objetivos}
% ----------------------------------------------------------

\subsection{Objetivo general}

Desarrollar una plataforma web con un agente autónomo capaz de percibir nuevos CV y postulaciones, analizar la compatibilidad con vacantes, generar resultados explicables y ejecutar acciones de soporte al reclutamiento con supervisión de RRHH.

\subsection{Objetivos específicos}

\begin{enumerate}[leftmargin=*]
    \item Implementar autenticación y autorización diferenciada para administradores de RRHH y candidatos.
    \item Permitir la gestión de vacantes con requisitos, pesos y umbrales configurables.
    \item Extraer texto y datos relevantes desde archivos PDF y DOCX.
    \item Calcular una puntuación ponderada basada en habilidades, experiencia y estudios.
    \item Separar el motor autónomo de la interfaz web mediante un worker independiente.
    \item Persistir eventos, reintentos y estados de procesamiento en SQLite.
    \item Generar rankings, justificaciones, reportes PDF y Excel.
    \item Preparar y enviar invitaciones reales mediante SMTP con aprobación previa de RRHH.
    \item Mantener trazabilidad mediante historial y monitoreo del worker.
\end{enumerate}

% ----------------------------------------------------------
\section{Alcance del sistema}
% ----------------------------------------------------------

\subsection{Funcionalidades para RRHH}

El administrador puede:

\begin{itemize}[leftmargin=*]
    \item visualizar un dashboard ejecutivo;
    \item registrar y cambiar el estado de vacantes;
    \item consultar candidatos;
    \item importar CV externos de forma masiva;
    \item solicitar reprocesos;
    \item visualizar el ranking y la decisión explicada del agente;
    \item aprobar borradores y enviar correos reales por SMTP;
    \item generar reportes;
    \item consultar el historial;
    \item monitorear el worker y la cola de eventos.
\end{itemize}

\subsection{Funcionalidades para candidatos}

El candidato puede:

\begin{itemize}[leftmargin=*]
    \item crear una cuenta e iniciar sesión;
    \item editar su perfil;
    \item cargar o actualizar su CV;
    \item consultar vacantes activas;
    \item postular a una vacante;
    \item visualizar sus postulaciones;
    \item revisar y responder invitaciones a entrevistas.
\end{itemize}

\subsection{Límites actuales}

\begin{itemize}[leftmargin=*]
    \item Los PDF escaneados sin capa de texto requieren incorporar OCR.
    \item La entrega final de correos depende de un servidor SMTP válido.
    \item La aplicación utiliza el servidor de desarrollo de Flask; para producción debe emplearse un servidor WSGI.
    \item El motor de análisis se basa en reglas explicables y un catálogo de habilidades; no utiliza un modelo generativo.
\end{itemize}

% ----------------------------------------------------------
\section{Arquitectura de la solución}
% ----------------------------------------------------------

\subsection{Vista general}

La arquitectura separa la interacción web del procesamiento autónomo:

\begin{figure}[ht]
\centering
\begin{tikzpicture}[
    node distance=1.35cm and 1.55cm,
    box/.style={rectangle, rounded corners, draw=primary, thick, fill=lightblue, align=center, minimum width=3.2cm, minimum height=1cm},
    db/.style={cylinder, shape border rotate=90, aspect=0.28, draw=secondary, thick, fill=lightblue, align=center, minimum width=2.8cm, minimum height=1.25cm},
    arrow/.style={-{Latex[length=2.5mm]}, thick, draw=darkgray}
]
    \node[box] (users) {Usuarios\\RRHH y candidatos};
    \node[box, right=of users] (flask) {Aplicación web\\Flask};
    \node[db, right=of flask] (sqlite) {SQLite\\datos y eventos};
    \node[box, below=of sqlite] (worker) {Worker autónomo\\\texttt{agent\_worker.py}};
    \node[box, left=of worker] (analysis) {Motor de análisis\\y ranking};
    \node[box, below=of analysis] (actions) {Reportes PDF/Excel\\y borradores SMTP};

    \draw[arrow] (users) -- (flask);
    \draw[arrow] (flask) -- (sqlite);
    \draw[arrow] (sqlite) -- (worker);
    \draw[arrow] (worker) -- (analysis);
    \draw[arrow] (analysis) -- (actions);
    \draw[arrow] (actions.east) -| (sqlite.south);
    \draw[arrow] (sqlite) -- (flask);
\end{tikzpicture}
\caption{Arquitectura general de \projectname{}.}
\label{fig:arquitectura}
\end{figure}

\subsection{Componentes principales}

\begin{longtable}{>{\raggedright\arraybackslash}p{3.5cm}>{\raggedright\arraybackslash}p{10.7cm}}
\caption{Componentes de software del proyecto.}\\
\toprule
\textbf{Componente} & \textbf{Responsabilidad} \\
\midrule
\endfirsthead
\toprule
\textbf{Componente} & \textbf{Responsabilidad} \\
\midrule
\endhead
\texttt{app.py} & Define rutas Flask, sesiones, control de acceso por rol, protección CSRF, carga de archivos y renderizado de plantillas. \\
\texttt{agent\_worker.py} & Ejecuta el loop autónomo, detecta CV nuevos en \texttt{uploads/}, reclama eventos persistentes y despacha tareas. \\
\texttt{modules/database.py} & Inicializa SQLite, realiza migraciones simples y concentra las operaciones de persistencia. \\
\texttt{modules/agent.py} & Orquesta ciclos masivos, postulaciones individuales, preanálisis de perfil y acciones posteriores. \\
\texttt{modules/extractor.py} & Extrae texto desde PDF o DOCX y reconoce correo, teléfono y nombre. \\
\texttt{modules/analyzer.py} & Detecta habilidades, experiencia y estudios; calcula compatibilidad; genera evidencia y justificación. \\
\texttt{modules/ranking.py} & Clasifica perfiles según umbrales configurables. \\
\texttt{modules/reports.py} & Exporta rankings en Excel y PDF. \\
\texttt{modules/email\_sender.py} & Gestiona la configuración SMTP y envía borradores aprobados por RRHH. \\
\texttt{templates/} & Contiene interfaces públicas, administrativas y de candidatos. \\
\bottomrule
\end{longtable}

% ----------------------------------------------------------
\section{Diseño del agente autónomo}
% ----------------------------------------------------------

\subsection{Ciclo operativo}

El agente implementa un ciclo percepción--razonamiento--decisión--acción--retroalimentación:

\begin{enumerate}[leftmargin=*]
    \item \textbf{Percibe:} recibe eventos por acciones web y detecta PDF o DOCX nuevos en \texttt{uploads/}.
    \item \textbf{Razona:} extrae texto, identifica datos y compara el CV contra la vacante.
    \item \textbf{Decide:} calcula compatibilidad y asigna un estado.
    \item \textbf{Actúa:} actualiza ranking, genera reportes y prepara invitaciones.
    \item \textbf{Retroalimenta:} registra ejecuciones, historial, errores y estado del worker.
\end{enumerate}

\subsection{Tipos de eventos}

\begin{table}[ht]
\centering
\begin{tabularx}{\textwidth}{>{\ttfamily}p{4.2cm}X}
\toprule
\textbf{Evento} & \textbf{Descripción} \\
\midrule
PREANALIZAR\_PERFIL & Compara el CV personal del candidato con todas las vacantes activas. \\
ANALIZAR\_POSTULACION & Analiza el CV de un candidato para una vacante elegida. \\
ANALIZAR\_IMPORTADOS & Procesa CV externos importados por RRHH o detectados en la carpeta de carga. \\
REPROCESAR\_VACANTE & Recalcula resultados para una vacante solicitada o reactivada. \\
\bottomrule
\end{tabularx}
\caption{Eventos procesados por el worker.}
\end{table}

\subsection{Persistencia, reintentos y recuperación}

La tabla \texttt{eventos\_agente} almacena el tipo de evento, su payload JSON, estado, número de intentos, límite máximo, disponibilidad temporal, worker responsable, bloqueo temporal y último error. El worker utiliza una transacción SQLite con \texttt{BEGIN IMMEDIATE} para reclamar un solo evento pendiente.

Cuando un worker reclama un evento, establece un arrendamiento temporal (\emph{lease}). Si el proceso se interrumpe antes de completar la tarea, otro ciclo recupera los eventos cuyo bloqueo haya expirado. Ante un error, el evento vuelve a estado pendiente hasta alcanzar el máximo de intentos. Esta estrategia evita perder trabajo cuando el proceso se reinicia.

\subsection{Monitoreo}

El worker actualiza periódicamente la tabla \texttt{estado\_worker}. El dashboard administrativo muestra:

\begin{itemize}[leftmargin=*]
    \item estado del proceso;
    \item último latido;
    \item cantidad de eventos pendientes, procesando y con error;
    \item lista reciente de eventos;
    \item ejecuciones autónomas recientes.
\end{itemize}

% ----------------------------------------------------------
\section{Modelo de análisis y compatibilidad}
% ----------------------------------------------------------

\subsection{Extracción de información}

El sistema acepta archivos PDF y DOCX. Para PDF utiliza \texttt{PyPDF2}; para DOCX utiliza \texttt{python-docx}. Después de obtener el texto, aplica expresiones regulares y reglas para extraer correo, teléfono, nombre, años de experiencia y nivel de estudios.

La detección de habilidades emplea coincidencias exactas delimitadas. Esto evita falsos positivos como detectar ``Java'' dentro de ``JavaScript''. El texto se normaliza a minúsculas, se eliminan acentos y se compactan espacios.

\subsection{Puntuación ponderada}

Cada vacante configura tres pesos:

\begin{itemize}[leftmargin=*]
    \item $w_h$: peso de habilidades;
    \item $w_e$: peso de experiencia;
    \item $w_s$: peso de estudios.
\end{itemize}

Los pesos se normalizan para que su suma sea uno:

\begin{equation}
\widehat{w_i} = \frac{w_i}{w_h + w_e + w_s}
\end{equation}

La puntuación de habilidades se calcula como:

\begin{equation}
P_h = \frac{\text{habilidades requeridas encontradas}}{\text{habilidades requeridas}} \times 100
\end{equation}

La puntuación de experiencia se limita a 100:

\begin{equation}
P_e =
\begin{cases}
100, & \text{si no se solicita experiencia} \\
\min\left(\dfrac{\text{años detectados}}{\text{años requeridos}} \times 100, 100\right), & \text{en otro caso}
\end{cases}
\end{equation}

La puntuación de estudios es 100 si el nivel detectado cumple o supera el nivel solicitado; de lo contrario es 0. Finalmente:

\begin{equation}
\text{Compatibilidad} =
P_h \widehat{w_h} +
P_e \widehat{w_e} +
P_s \widehat{w_s}
\end{equation}

\subsection{Clasificación}

Cada vacante define un umbral de preselección $U_p$ y un umbral de observación $U_o$:

\begin{equation}
\text{Estado} =
\begin{cases}
\text{Preseleccionado}, & \text{si Compatibilidad} \ge U_p \\
\text{En observación}, & \text{si Compatibilidad} \ge U_o \\
\text{No recomendado}, & \text{en otro caso}
\end{cases}
\end{equation}

La configuración predeterminada asigna 70\% a habilidades, 20\% a experiencia y 10\% a estudios. Los umbrales predeterminados son 80\% para preselección y 60\% para observación.

\subsection{Explicabilidad}

El resultado almacena:

\begin{itemize}[leftmargin=*]
    \item habilidades encontradas y faltantes;
    \item experiencia detectada;
    \item estudios detectados;
    \item desglose de criterios;
    \item justificación textual;
    \item acción sugerida.
\end{itemize}

Por tanto, RRHH puede revisar el fundamento de cada resultado y no solo un número aislado.

% ----------------------------------------------------------
\section{Modelo de datos}
% ----------------------------------------------------------

\begin{longtable}{>{\ttfamily}p{3.3cm}>{\raggedright\arraybackslash}p{10.9cm}}
\caption{Tablas principales de SQLite.}\\
\toprule
\textbf{Tabla} & \textbf{Propósito} \\
\midrule
\endfirsthead
\toprule
\textbf{Tabla} & \textbf{Propósito} \\
\midrule
\endhead
usuarios & Cuentas, credenciales hash, rol, perfil y referencia al CV personal. \\
vacantes & Información de puestos, requisitos, pesos, umbrales, entrevista y estado. \\
candidatos & Perfiles analizados, incluidos CV externos y perfiles vinculados a usuarios. \\
analisis & Compatibilidad, clasificación, evidencia y explicación por candidato y vacante. \\
postulaciones & Relación entre candidatos registrados y vacantes elegidas. \\
correos & Borradores e invitaciones enviadas, asociados a candidato y vacante. \\
historial & Bitácora funcional del sistema. \\
ejecuciones\_agente & Resumen de cada ciclo autónomo, origen, estado y resultado. \\
eventos\_agente & Cola persistente, reintentos, lease y errores del worker. \\
estado\_worker & Último latido y estado visible del worker. \\
\bottomrule
\end{longtable}

% ----------------------------------------------------------
\section{Seguridad}
% ----------------------------------------------------------

\subsection{Separación de roles}

El sistema distingue:

\begin{itemize}[leftmargin=*]
    \item \textbf{Administrador/RRHH:} accede a rutas administrativas y datos globales.
    \item \textbf{Candidato:} accede únicamente a su perfil, CV, postulaciones e invitaciones.
    \item \textbf{Agente:} no es un usuario; funciona como proceso interno.
\end{itemize}

Los decoradores \texttt{role\_required} y \texttt{login\_required} validan la sesión antes de ejecutar vistas protegidas. Cuando un usuario intenta ingresar a un área ajena a su rol, la aplicación lo redirige a su dashboard correspondiente.

\subsection{Medidas aplicadas}

\begin{itemize}[leftmargin=*]
    \item Contraseñas almacenadas mediante hash de Werkzeug.
    \item Tokens CSRF para solicitudes POST.
    \item Límite de carga de 16 MB.
    \item Nombres de archivos saneados mediante \texttt{secure\_filename}.
    \item Validación de extensión y firma básica del archivo.
    \item Consultas SQL parametrizadas.
    \item Variables de entorno para secretos y credenciales SMTP.
    \item Botón SMTP deshabilitado si falta configuración.
    \item Conservación del borrador cuando un envío falla.
\end{itemize}

\subsection{Configuración SMTP}

El envío real se habilita con variables de entorno:

\begin{lstlisting}[language=bash,caption={Ejemplo de configuración SMTP para Gmail en PowerShell.}]
$env:SMTP_HOST="smtp.gmail.com"
$env:SMTP_PORT="587"
$env:SMTP_USER="tu-cuenta@gmail.com"
$env:SMTP_PASSWORD="tu-contrasena-de-aplicacion"
$env:SMTP_FROM="tu-cuenta@gmail.com"
$env:SMTP_USE_TLS="1"
python app.py
\end{lstlisting}

Para Gmail debe emplearse una contraseña de aplicación. Las credenciales no deben escribirse en archivos versionados ni compartirse en el código fuente.

% ----------------------------------------------------------
\section{Flujos de uso}
% ----------------------------------------------------------

\subsection{Carga de CV personal}

\begin{enumerate}[leftmargin=*]
    \item El candidato inicia sesión y carga un PDF o DOCX.
    \item Flask valida y guarda el archivo.
    \item La base de datos crea un evento \texttt{PREANALIZAR\_PERFIL}.
    \item El worker reclama el evento.
    \item El agente compara el CV con vacantes activas.
    \item Se guardan recomendaciones y evidencia.
    \item Si existen postulaciones previas, se encolan análisis específicos.
\end{enumerate}

\subsection{Postulación}

\begin{enumerate}[leftmargin=*]
    \item El candidato elige una vacante activa.
    \item La aplicación registra la postulación.
    \item Se encola \texttt{ANALIZAR\_POSTULACION}.
    \item El agente procesa el CV, actualiza la compatibilidad y clasifica el perfil.
    \item El ranking y los reportes se actualizan.
\end{enumerate}

\subsection{Importación externa}

\begin{enumerate}[leftmargin=*]
    \item RRHH importa varios CV o deposita archivos en \texttt{uploads/}.
    \item El worker detecta archivos nuevos no registrados.
    \item Se crea un evento \texttt{ANALIZAR\_IMPORTADOS} por vacante activa.
    \item El agente procesa los CV y actualiza los rankings.
\end{enumerate}

\subsection{Envío de invitaciones}

\begin{enumerate}[leftmargin=*]
    \item El agente prepara borradores para perfiles preseleccionados.
    \item RRHH revisa la pantalla de correos.
    \item RRHH aprueba el envío.
    \item El módulo SMTP envía cada mensaje individualmente.
    \item Los mensajes exitosos cambian a estado \texttt{Enviado}.
    \item Los fallidos permanecen como borradores y se registran en el historial.
\end{enumerate}

% ----------------------------------------------------------
\section{Tecnologías utilizadas}
% ----------------------------------------------------------

\begin{table}[ht]
\centering
\begin{tabularx}{\textwidth}{>{\bfseries}p{3.6cm}X}
\toprule
Tecnología & Uso en el proyecto \\
\midrule
Python & Lenguaje principal. \\
Flask & Aplicación web, rutas, sesiones y plantillas. \\
SQLite & Persistencia local y cola de eventos. \\
Jinja2 & Renderizado de interfaces HTML. \\
Bootstrap & Diseño visual responsivo con recursos locales. \\
PyPDF2 & Extracción de texto desde PDF. \\
python-docx & Lectura de documentos DOCX. \\
pandas y openpyxl & Exportación de reportes Excel. \\
FPDF & Generación de reportes PDF. \\
smtplib & Envío real de invitaciones mediante SMTP. \\
unittest & Pruebas automatizadas. \\
\bottomrule
\end{tabularx}
\caption{Tecnologías empleadas.}
\end{table}

% ----------------------------------------------------------
\section{Instalación y ejecución}
% ----------------------------------------------------------

\subsection{Preparación}

\begin{lstlisting}[language=bash,caption={Instalación del entorno.}]
python -m venv .venv
.venv\Scripts\activate
pip install -r requirements.txt
\end{lstlisting}

\subsection{Ejecución}

La arquitectura requiere dos procesos:

\begin{lstlisting}[language=bash,caption={Terminal 1: aplicación web.}]
python app.py
\end{lstlisting}

\begin{lstlisting}[language=bash,caption={Terminal 2: worker autónomo.}]
python agent_worker.py
\end{lstlisting}

La interfaz queda disponible en:

\begin{center}
\url{http://127.0.0.1:5000/}
\end{center}

La cuenta administrativa de demostración es:

\begin{center}
\texttt{admin@demo.com} / \texttt{admin123}
\end{center}

Para un despliegue real debe cambiarse la clave secreta y eliminarse o modificarse la cuenta demo.

% ----------------------------------------------------------
\section{Pruebas y validación}
% ----------------------------------------------------------

\subsection{Estrategia}

El proyecto incluye pruebas automatizadas ejecutables mediante:

\begin{lstlisting}[language=bash]
python -m unittest discover -s tests -v
\end{lstlisting}

Durante la validación se ejecutaron 14 pruebas con resultado satisfactorio. Además, se verificó la compilación de módulos Python mediante \texttt{py\_compile}.

\subsection{Casos cubiertos}

\begin{longtable}{>{\raggedright\arraybackslash}p{8.2cm}>{\centering\arraybackslash}p{2.2cm}>{\raggedright\arraybackslash}p{3.5cm}}
\caption{Resumen de validaciones automatizadas.}\\
\toprule
\textbf{Caso} & \textbf{Resultado} & \textbf{Objetivo} \\
\midrule
\endfirsthead
\toprule
\textbf{Caso} & \textbf{Resultado} & \textbf{Objetivo} \\
\midrule
\endhead
Login administrativo y bloqueo de rutas para candidatos & Correcto & Control de acceso \\
Protección CSRF en solicitudes POST & Correcto & Seguridad web \\
Rechazo de vacante inexistente & Correcto & Integridad funcional \\
Rechazo de PDF falso por extensión & Correcto & Validación de archivos \\
Detección exacta de habilidades & Correcto & Calidad del análisis \\
Cálculo ponderado de experiencia y estudios & Correcto & Compatibilidad \\
Ocultamiento de vacantes cerradas & Correcto & Reglas del negocio \\
Umbrales configurables & Correcto & Clasificación \\
Preanálisis sin postulación automática & Correcto & Separación de acciones \\
Invitaciones sin duplicados & Correcto & Integridad de correos \\
Envío limitado a la vacante activa & Correcto & Alcance SMTP \\
Procesamiento desde cola persistente & Correcto & Autonomía \\
Error tras agotar reintentos & Correcto & Tolerancia a fallos \\
Detección de CV copiado directamente a \texttt{uploads/} & Correcto & Percepción autónoma \\
\bottomrule
\end{longtable}

\subsection{Validación SMTP}

Las pruebas de correo reemplazan temporalmente el transporte SMTP por un mock. Esto verifica que:

\begin{itemize}[leftmargin=*]
    \item el destinatario, asunto y mensaje se entregan al módulo de transporte;
    \item solo se procesan borradores de la vacante elegida;
    \item un envío exitoso cambia a estado \texttt{Enviado};
    \item otros borradores permanecen intactos.
\end{itemize}

No se incluyen credenciales reales en el repositorio.

% ----------------------------------------------------------
\section{Resultados obtenidos}
% ----------------------------------------------------------

La solución implementada cumple con los objetivos planteados:

\begin{itemize}[leftmargin=*]
    \item existe separación funcional entre RRHH y candidatos;
    \item el agente funciona como proceso interno y no como usuario;
    \item el análisis se ejecuta fuera de las solicitudes HTTP;
    \item la cola sobrevive reinicios;
    \item el worker detecta nuevos archivos y eventos;
    \item los resultados incluyen evidencia y justificación;
    \item se generan reportes PDF y Excel;
    \item los correos pueden enviarse realmente mediante SMTP;
    \item RRHH dispone de trazabilidad y monitoreo.
\end{itemize}

% ----------------------------------------------------------
\section{Trabajo futuro}
% ----------------------------------------------------------

\begin{enumerate}[leftmargin=*]
    \item Incorporar OCR para PDF escaneados mediante Tesseract u otro servicio especializado.
    \item Migrar a PostgreSQL para despliegues con mayor concurrencia.
    \item Ejecutar Flask detrás de Gunicorn, Waitress o un servidor WSGI equivalente.
    \item Incorporar HTTPS y gestión formal de secretos.
    \item Añadir edición de plantillas de correo desde el panel.
    \item Implementar métricas de entrega SMTP y seguimiento de rebotes.
    \item Añadir pruebas de integración con un servidor SMTP de pruebas.
    \item Expandir el catálogo de habilidades o incorporar modelos NLP supervisados.
    \item Permitir programación horaria de ciclos autónomos.
    \item Incorporar mecanismos de auditoría y anonimización para cumplimiento normativo.
\end{enumerate}

% ----------------------------------------------------------
\section{Conclusiones}
% ----------------------------------------------------------

\projectname{} demuestra que es posible construir un agente autónomo de reclutamiento con herramientas accesibles y una arquitectura comprensible. La plataforma no se limita a presentar pantallas: observa eventos, recupera tareas persistentes, procesa CV en segundo plano y ejecuta acciones derivadas.

El enfoque de compatibilidad ponderada aporta explicabilidad. RRHH puede revisar habilidades, experiencia, estudios, faltantes y acciones sugeridas antes de tomar decisiones. La automatización no reemplaza la supervisión humana: la fortalece al reducir tareas repetitivas y mantener trazabilidad.

Finalmente, la separación entre Flask y el worker mejora la robustez conceptual del proyecto. El agente continúa trabajando aunque no exista una interacción web inmediata, recupera eventos interrumpidos y reporta su estado. Esta separación convierte la propuesta en una solución autónoma funcional y extensible.

% ----------------------------------------------------------
\section*{Referencias}
\addcontentsline{toc}{section}{Referencias}
% ----------------------------------------------------------

\begin{enumerate}[leftmargin=*]
    \item Pallets Projects. \textit{Flask Documentation}. Disponible en: \url{https://flask.palletsprojects.com/}
    \item Python Software Foundation. \textit{sqlite3 --- DB-API 2.0 interface for SQLite databases}. Disponible en: \url{https://docs.python.org/3/library/sqlite3.html}
    \item Python Software Foundation. \textit{smtplib --- SMTP protocol client}. Disponible en: \url{https://docs.python.org/3/library/smtplib.html}
    \item SQLite Consortium. \textit{SQLite Documentation}. Disponible en: \url{https://www.sqlite.org/docs.html}
    \item Bootstrap Team. \textit{Bootstrap Documentation}. Disponible en: \url{https://getbootstrap.com/docs/}
    \item PyPDF2 Contributors. \textit{PyPDF2 Documentation}. Disponible en: \url{https://pypdf2.readthedocs.io/}
    \item python-docx Contributors. \textit{python-docx Documentation}. Disponible en: \url{https://python-docx.readthedocs.io/}
\end{enumerate}

% ----------------------------------------------------------
\appendix
\section{Estructura del proyecto}
% ----------------------------------------------------------

\begin{lstlisting}[caption={Estructura principal de carpetas.}]
app.py                  Aplicacion web Flask
agent_worker.py         Worker autonomo
modules/
  agent.py              Orquestacion
  analyzer.py           Compatibilidad explicable
  database.py           Persistencia SQLite
  email_sender.py       Envio SMTP
  extractor.py          Extraccion PDF y DOCX
  ranking.py            Clasificacion
  reports.py            Reportes PDF y Excel
templates/              Interfaces HTML
static/                 CSS, JavaScript y Bootstrap local
uploads/                CV cargados
reports/                Reportes generados
database/               Base SQLite
logs/                   Registro del worker
tests/                  Pruebas automatizadas
\end{lstlisting}

\section{Variables de entorno}

\begin{longtable}{>{\ttfamily}p{4.1cm}>{\raggedright\arraybackslash}p{10.1cm}}
\toprule
\textbf{Variable} & \textbf{Uso} \\
\midrule
SECRET\_KEY & Clave privada de sesión Flask. \\
SMTP\_HOST & Servidor SMTP, por ejemplo \texttt{smtp.gmail.com}. \\
SMTP\_PORT & Puerto SMTP; el valor habitual con TLS es \texttt{587}. \\
SMTP\_USER & Usuario utilizado para autenticarse. \\
SMTP\_PASSWORD & Contraseña de aplicación o secreto SMTP. \\
SMTP\_FROM & Dirección remitente; si se omite se usa \texttt{SMTP\_USER}. \\
SMTP\_USE\_TLS & Activa STARTTLS cuando vale \texttt{1}. \\
SMTP\_USE\_SSL & Activa conexión SSL directa cuando vale \texttt{1}. \\
\bottomrule
\end{longtable}

\section{Comando de validación}

\begin{lstlisting}[language=bash]
python -m py_compile app.py agent_worker.py modules/database.py \
    modules/email_sender.py
python -m unittest discover -s tests -v
\end{lstlisting}

\end{document}
